\chapter{Literature Review}
\label{chap:related-work}

\section{Scope and comparison perspective}

Remote sensing scene classification connects visual representation learning
with a question about spatial organisation: which arrangements of land-cover
elements distinguish one scene from another? Cheng et al.~\cite{cheng2020scene}
reviewed this problem through its methods, benchmarks and remaining challenges.
Thapa et al.~\cite{thapa2023review} subsequently examined deep-learning-based
scene classification. These reviews provide a broad context for the present
study, whose focus is the narrower relationship between contextual image
features, regional aggregation and explicit graph reasoning.

This chapter follows three connected streams: visual backbones for remote
sensing, graph-based image understanding, and region grouping. It compares
methods along five dimensions: the visual unit represented by a node or token;
the rule defining its neighbours; the operation aggregating their information;
the position of the relational module; and the supervision used to learn the
representation. This organisation separates architectural similarity from
evidence that two methods solve the same task under comparable conditions.

Scene classification assigns an image-level description. A single-label
benchmark selects one scene category, whereas a multi-label benchmark permits
several land-cover labels for the same image. Neither task directly requires
object boxes or semantic masks. Thus, a region useful for classification need
not coincide with an annotated object. This distinction matters when grouping
and graph visualisations are interpreted: class supervision can reward a
predictive representation without enforcing a human-defined segmentation.

The chapter includes general vision methods where their mechanisms directly
inform GAViT. Their inclusion does not imply that their natural-image,
video, point-cloud or segmentation results transfer unchanged to remote
sensing classification. Likewise, the review distinguishes a design precedent
from a reproduced experimental comparator.

\section{Learning contextual visual representations}

\subsection{Convolutional features and transfer to remote sensing}

Convolutional neural networks (CNNs) learn spatially shared feature extractors
whose receptive fields expand through successive layers. Simonyan and
Zisserman~\cite{simonyan2015vgg} investigated deep networks built from small
convolution filters. He et al.~\cite{he2016resnet} introduced residual learning
to facilitate optimisation of substantially deeper networks. These designs
provide two relevant foundations: hierarchical features describe images at
different levels, and residual paths allow an additional transformation to
modify an existing representation. Neither property alone specifies an
explicit graph between image regions.

The usefulness of natural-image pretraining for aerial scenes was investigated
directly. Penatti et al.~\cite{penatti2015deepfeatures} examined whether deep
features generalised from everyday objects to remote sensing and aerial
scenes. Hu et al.~\cite{hu2015transferring} studied transferring pretrained
CNN representations to high-resolution scene classification. Nogueira
et al.~\cite{nogueira2017exploiting} further investigated ways of exploiting
CNNs for this task. Together, these studies motivate treating pretraining and
adaptation as experimental factors, rather than attributing every performance
difference to the last architectural module.

Scale handling is another example of a distinct design choice. Xie
et al.~\cite{xie2019scalefree} investigated a scale-free CNN that replaced
fully connected processing with convolutional operations and global pooling
to accommodate variable image sizes. Accommodating scale variation through
the feature extractor differs from constructing content-adaptive region
membership. These mechanisms may address related scene variation without
being equivalent architectural interventions.

Transfer also affects the interpretation of a common-backbone comparison.
Using the same pretrained feature extractor provides a controlled starting
point, but joint fine-tuning can change that extractor differently in two
models. A graph-augmented model therefore includes both its added relational
operations and the representation learned in their presence. Removing its
graph after training does not recover an independently trained backbone
baseline. This is an experimental distinction, rather than a limitation
specific to CNNs or Transformers.

CNNs remain a relevant comparison class. Liu et al.~\cite{liu2022convnext}
revisited convolutional design and introduced ConvNeXt through a systematic
modernisation of a residual network. This work cautions against explaining
architectural performance solely by a CNN-versus-Transformer label. For the
present study, the main backbone is held fixed so that the regional branch
can be examined without simultaneously replacing the entire feature hierarchy.

\subsection{Attention, image tokens and hierarchical backbones}

Vaswani et al.~\cite{vaswani2017attention} introduced the Transformer around
attention-based sequence processing. Its relevance here is the content-dependent
combination of representations, rather than the original language task.
Dosovitskiy et al.~\cite{dosovitskiy2021vit} applied a Transformer to image
patch sequences in the Vision Transformer (ViT). Patch tokenisation supplied
a regular input representation; attention subsequently contextualised those
tokens. The resulting features should not be interpreted as independent
unprocessed image blocks.

Training choices remained consequential. Touvron
et al.~\cite{touvron2021deit} studied data-efficient image Transformers and
attention-based distillation. Their work demonstrated that architecture,
training recipe and teacher information must be considered together. A
comparison using different pretraining data or distillation supervision does
not isolate the effect of a regional graph, even when both models contain
Transformer layers.

Hierarchical Transformers provided another route to organising spatial
information. Wang et al.~\cite{wang2021pvt} proposed the Pyramid Vision
Transformer, which progressively reduced spatial resolution to form a
feature hierarchy. Liu et al.~\cite{liu2021swin} used local windows and
shifted-window connections in Swin. These approaches made multi-scale and
cross-location interaction part of the backbone design. Consequently, adding
a graph above Swin tests a further relational organisation of features that
already contain context; it does not introduce context into an otherwise
context-free model.

Backbone token resolution constrains the regional branch. A late-stage
representation is computationally compact and carries high-level information,
but its spatial units can mix several land-cover elements. Moving a graph to
earlier features would change its node count, receptive fields and cost.
Similarly, treating every patch as a node differs from aggregating several
contextual tokens into one region. This motivates examining node construction
and insertion depth alongside the choice of graph operator.

\subsection{Remote sensing Transformers and domain pretraining}

Bazi et al.~\cite{bazi2021remotevit} investigated vision Transformers for
remote sensing image classification. Zhang et al.~\cite{zhang2021trs}
also studied Transformer-based remote sensing scene representation. Their
relevance is the application of token interactions to scene-level recognition;
using attention does not by itself imply that the model learns explicit
geographic objects or their semantic relationships.

Xu et al.~\cite{xu2025hetmcl} more recently combined high-frequency-enhancing
Transformer processing with multi-layer contextual learning for remote
sensing scene classification. This provides a further example of treating
feature detail and contextual fusion jointly. Its frequency-based design
differs from defining a small graph of pooled regions, and its reported
results do not identify which mechanism would improve the shared-backbone
protocol used here.

Recent representation learning also changes what is available before task
fine-tuning. He et al.~\cite{he2022mae} proposed masked autoencoders, using
an asymmetric encoder--decoder to reconstruct missing image content. In
remote sensing, Cong et al.~\cite{cong2022satmae} considered temporal and
multispectral information in masked-autoencoder pretraining. Reed
et al.~\cite{reed2023scalemae} incorporated scale awareness, while Noman
et al.~\cite{noman2024satmaepp} investigated multi-scale reconstruction.
These works make the pretraining signal and spatial scale important parts
of the representation-learning problem.

Liu et al.~\cite{liu2024remoteclip} developed RemoteCLIP through visual--language
representation learning for remote sensing. Such supervision differs from
learning only from the image-level labels of a downstream scene dataset.
Additional text, temporal observations or spectral bands can improve the
information available to a model, but they also change the comparison
conditions. This thesis uses RGB inputs and a common recorded pretrained
backbone in its internal architectural comparisons. Domain-pretrained methods
are therefore relevant context, rather than automatically interchangeable
baselines under that protocol.

Across these studies, richer pretraining and explicit relational structure
address different design choices. Stronger initial features may reduce the
remaining benefit of an added module, or may make meaningful grouping easier.
Neither possibility follows from architectural inspection alone. Holding
the backbone initialisation and downstream data protocol fixed is useful for
testing one specific regional design, while broader backbone comparisons
would answer a separate question about representation quality.

\section{Explicit graphs for visual understanding}

\subsection{Nodes, adjacency and message aggregation}

A graph representation requires a definition of its entities and connections.
For an image, nodes may describe pixels, patches, pooled regions or objects.
Spatial adjacency, feature similarity and known semantic relations encode
different assumptions about which nodes should exchange information. A
network that predicts scene labels from a feature-neighbour graph should
therefore be distinguished from one that predicts labelled subject--relation--object
triples in a conventional scene graph.

Kipf and Welling~\cite{kipf2017gcn} developed a graph convolutional network
(GCN) using normalised neighbourhood propagation. Hamilton
et al.~\cite{hamilton2017graphsage} proposed learning neighbourhood
aggregation functions for inductive node representation. Veli\v{c}kovi\'{c}
et al.~\cite{velickovic2018gat} introduced graph attention to learn how
neighbours contribute within a supplied connectivity structure. These works
established useful aggregation mechanisms, but did not establish that a
particular image grouping or feature-distance rule is semantically correct.

The distinction between topology and aggregation is central to GAViT.
Selecting a neighbour by cosine similarity answers which representation is
eligible to contribute. Attention then determines its contribution according
to the graph layer. The two scores need not be identical. A fixed spatial
edge can receive a learned attention weight, while a content-selected edge
can be used without passing its original similarity value as a message
coefficient. Reporting a graph therefore requires specifying both decisions.

Directed edges add another necessary detail. If a query selects a feature
neighbour, the intended aggregation usually requires the neighbour's
representation to reach the query. Reversing the stored direction can change
that operation. The explicit direction, treatment of self-connections and
duplicate spatial/feature edges are consequently part of the method
definition, even when they are not visible in a schematic diagram.

\subsection{Dynamic feature graphs and visual backbones}

Wang et al.~\cite{wang2019dgcnn} introduced EdgeConv in a dynamic graph
network for point clouds. Its feature-space neighbourhoods provide a useful
precedent for content-dependent connectivity. However, a point-cloud graph
and a regional image graph have different geometric inputs and prediction
settings. The precedent supports a design option, not a direct transfer of
the reported task performance.

Han et al.~\cite{han2022vig} represented images through patch nodes in
Vision GNN (ViG). Their design combined a graph aggregation module with
feature transformations in a visual backbone. Munir
et al.~\cite{munir2024greedyvig} subsequently investigated efficient graph
construction for vision graph networks. These works connect connectivity
choices to computational behaviour: the cost of constructing neighbours is
part of an image graph model, rather than a free preprocessing assumption.

Graph recomputation is also a separate architectural choice. A model may
construct a feature graph once and reuse it through several message-passing
layers, or update neighbours after intermediate features change. Both may
be described as content-dependent, but their information paths and costs
differ. GAViT uses a low-node-count regional branch above a pretrained visual
hierarchy, which differs from replacing that hierarchy with repeated graph
blocks. Its evaluation must include the pooling, graph construction and
feedback that create the branch, not only the graph layer in isolation.

\subsection{Graph and attention as complementary operators}

Chen et al.~\cite{chen2019glore} proposed Graph-Based Global Reasoning
Networks, which projected coordinate-space features to a smaller interaction
space using weighted pooling, performed graph convolution, and broadcast
the resulting information back. This is a particularly direct precedent
for aggregation, graph reasoning and return to a spatial representation.
The distinction of GAViT therefore concerns its anchored grouping, sparse
spatial-and-feature topology and late-stage feedback implementation, rather
than the general existence of this three-step organisation.

Graph networks and Transformers can be combined at different levels.
Ramp\'{a}\v{s}ek et al.~\cite{rampasek2022graphgps} studied a framework that
combined local message passing, global attention and structural information
for graph learning. This offered a general perspective on complementary
information paths, although its graph benchmarks were not a direct evaluation
of regional remote sensing image classification.

For visual models, complementary paths can appear as parallel feature
branches, alternating blocks or late-stage augmentation. Their meanings
depend on the units exchanged between the branches. A channel-attention
operation, token self-attention and a graph over pooled regions are not
interchangeable merely because all use learned weights. Comparing their
insertion points and node granularity is more informative than grouping
them under a single attention label.

This observation also bounds claims of novelty. Combining graph learning
with attention-based vision has substantial precedent. The contribution of
a particular system can instead concern a precisely specified construction,
its controlled comparison and what its analysis reveals about the conditions
under which the added structure is useful.

\section{From fixed partitions to learned token groups}

\subsection{Fixed grids and hard feature grouping}

Fixed spatial partitions offer stable region identifiers and inexpensive
membership rules. They also provide an immediate interpretation for spatial
neighbours. Their boundaries are determined by coordinates, however, so a
linear structure or mixed scene may cross several groups. Learned weighting
within a fixed group changes its pooled representation without changing
which tokens belong to it. This distinction is particularly important when
the term attention-guided grouping is used.

Hard feature grouping offers a different trade-off. Clustering can place
similar tokens together even when their positions differ, but the resulting
groups may lose a consistent spatial identity. The choice of feature metric,
number of groups, initialisation and handling of empty groups then becomes
part of the model. A discrete membership rule also differs from a soft
assignment through which downstream gradients can change token membership.
These alternatives motivate retaining a fixed grouping as a reference
rather than assuming that increased flexibility must improve classification.

\subsection{Superpixels, locality and differentiable association}

Achanta et al.~\cite{achanta2012slic} proposed Simple Linear Iterative
Clustering (SLIC) superpixels using a local combination of appearance and
position. Jampani et al.~\cite{jampani2018ssn} introduced Superpixel Sampling
Networks (SSN), connecting differentiable superpixel formation with learned
features and downstream objectives. The shared lesson for this thesis is
that content dependence and spatial constraints can coexist. They need not
be treated as mutually exclusive grouping strategies.

The level at which grouping operates nevertheless matters. Pixel-level
superpixels can be assessed against fine spatial boundaries, whereas a
group of late-stage contextual tokens may have a much coarser and overlapping
relationship to the input. A fixed anchor can provide a useful prior without
being a true object centre. Conversely, a visually irregular group does not
automatically imply a poor classification representation. The appropriate
claim depends on the supervision, spatial resolution and measured outcome.

Differentiable association also requires specifying its normalisation. A
token may distribute membership mass across several regions, while a region
separately normalises the weights of all contributing tokens. These answer
different questions: where a token belongs, and how a region is represented.
Separating the two is useful when the same membership is subsequently used
to send region information back to the original tokens.

\subsection{Object-centric and semantically supervised groups}

Locatello et al.~\cite{locatello2020slot} introduced Slot Attention to
organise input features through a set of competing slots and iterative
updates. Its grouping mechanism provides a reference for differentiating
input membership from slot aggregation. Its exchangeable slot identities
also differ from regions bound to a persistent spatial anchor. A graph
defined by neighbouring anchor indices cannot be transferred to arbitrary
slots without explaining the new meaning of adjacency.

Xu et al.~\cite{xu2022groupvit} proposed GroupViT, in which semantic
segmentation emerged through image--text supervision and hierarchical visual
grouping. It illustrates how the learning signal and grouping architecture
interact. GAViT uses scene classification labels rather than the same
visual--language supervision, so it cannot inherit GroupViT's semantic
segmentation interpretation solely by using learned assignments.

These methods also show why soft grouping should be specified concretely.
The term can describe competitive slots, local cluster associations or
hierarchical group-token assignments. They need not use the same update
rule, locality, forward representation or gradient estimator. A negative
result for a single local-prototype design therefore does not reject
learnable grouping as a whole; a positive classification result would
likewise not establish object discovery without an appropriate evaluation.

\subsection{Token aggregation, interaction and reconstruction}

Wu et al.~\cite{wu2020visualtransformers} explored visual Transformers that
organised spatial features through visual tokens. Ryoo
et al.~\cite{ryoo2021tokenlearner} introduced TokenLearner and TokenFuser in
a video representation framework, combining adaptive token extraction with
a mechanism for returning information to spatial features. These provide
prior examples of reducing the number of units used for interaction while
retaining a connection to a denser visual representation.

Huang et al.~\cite{huang2023stvit} proposed a vision Transformer with super
token sampling, denoted STViT: sparse associations aggregated visual tokens,
interaction occurred among super tokens, and information was mapped back.
Chang et al.~\cite{chang2023semanticstvit} introduced Semantic Token ViT,
which uses the same STViT abbreviation for a different architecture. Its
STViT-R variant restores spatial representations through attention from
image-token queries to semantic-token keys and values. Both studies provide
prior examples of compact token interaction and return to a denser spatial
representation. GAViT uses its particular regional assignments and sparse
graph operator within this established organisational pattern.

Bolya et al.~\cite{bolya2023tome} proposed Token Merging to reduce token
redundancy through merging. This makes computational efficiency a distinct
purpose for grouping. A method that reduces tokens inside a backbone may
save subsequent computation, whereas a regional branch added after a full
backbone introduces its own cost. Similar reductions in node count therefore
do not imply similar overall speed or memory behaviour.

The local prototype variant in this thesis occupies a restricted position
among these alternatives. It retains anchor neighbourhoods, constructs
prototypes from each image's contextual tokens, and uses one local assignment
step before sparse graph reasoning. It neither implements the complete
iterative Slot Attention procedure nor replaces super-token attention with
an equivalent operator. Its value is tested as a specified grouping and
feedback variant, rather than inferred from the success of a related family.

\section{Directly related remote sensing graph methods}

\subsection{Hierarchical graph-enhanced Transformers}

Li et al.~\cite{li2024hgtnet} proposed a Hierarchical Graph-Enhanced
Transformer Network (HGTNet) for remote sensing scene classification. The
method combined hierarchical Transformer representations, attention modules
and graph-based feature extraction. It is directly relevant to the claim
that explicit relations can complement visual attention in this application.
Its placement and organisation of modules differ from a single regional
graph above the final stage of a shared Swin-T backbone.

This difference makes HGTNet a complete-method comparator rather than a
single-component ablation. A comparison must retain and disclose its
backbone, graph operations, feature fusion and pretraining assumptions.
Replacing those elements with an available Swin--GAT sequence would create
a different implementation. Its published performance also depends on its
reported benchmark protocol; it cannot be inserted into a common-protocol
result column as though it were measured in this project.

\subsection{Deep feature aggregation through graph convolution}

Xu et al.~\cite{xu2022dfagcn} proposed a Deep Feature Aggregation Framework
driven by a Graph Convolutional Network (DFAGCN). Their method combined
multiple levels of convolutional features, graph-based aggregation, feature
fusion and a final linear classification stage. Its relationship to GAViT
lies in aggregating visual information through an explicit graph, while
its feature source and processing chain differ from the token-feedback
design studied here.

The complete chain matters when defining a comparison. Multiscale feature
extraction, graph training, fusion coefficients and the classifier's
selection procedure can each change the resulting representation. A
standalone graph-convolution layer attached to the present classifier is
therefore not a reproduction of DFAGCN. Independent implementation choices
must be recorded where the available specification leaves uncertainty.
Preparing or verifying individual components also does not establish a
completed benchmark result.

Table~\ref{tab:relational-designs} summarises the architectural relationships
most relevant to the study. Its entries compare designs, not performance
under a shared experimental protocol.

\begin{table}[tbp]
 \centering
 \caption{Selected relational designs and their relationship to GAViT.
 The table describes architectural organisation rather than a ranking of
 published classification scores.}
 \label{tab:relational-designs}
 \small
 \setlength{\tabcolsep}{4pt}
 \begin{tabular}{p{2.3cm}p{4.0cm}p{6.6cm}}
  \toprule
  Method & Visual units and interaction & Relationship to the present study\\
  \midrule
  GloRe (2019) & Pooled interaction-space graph and broadcasting & Direct graph precedent for aggregation, relational reasoning and mapping back.\\[0.4em]
  ViG (2022) & Patch-node graph backbone & Repeated visual graph processing differs from a late regional branch.\\[0.4em]
  TokenLearner (2021) & Adaptive tokens and spatial fusion & Aggregation and return to spatial features predate the present framework.\\[0.4em]
  Super Token Sampling (2023) & Sparse token association and super-token attention & Shares the broad aggregation--interaction--mapping organisation; the interaction operator differs.\\[0.4em]
  HGTNet (2024) & Hierarchical Transformer and graph feature fusion & Direct remote sensing comparator with a different full architecture.\\[0.4em]
  DFAGCN (2022) & Multilevel CNN features and graph aggregation & Direct graph-based comparator with a different staged training and classification chain.\\[0.4em]
  GAViT & Anchored token groups and sparse regional GAT & Tests fixed and local-prototype membership with a common backbone and feedback interface.\\
  \bottomrule
 \end{tabular}
\end{table}

\section{Evaluation, interpretability and computational trade-offs}

Architectural comparisons require compatible evidence. The number of training
images, validation rule, pretraining, augmentation and numerical settings
can all affect the meaning of a result. Fixed-split training repetitions
measure a different source of variation from experiments that redraw the
data partition. Published means and standard deviations should therefore
be accompanied by their repetition protocol. A higher number in a paper
using another split or additional supervision is not an isolated estimate
of the benefit of its graph module.

Single-label and multi-label metrics also answer different questions. Overall
accuracy counts correct scene decisions, whereas average precision assesses
the ranking of examples for a label. Macro-averaged measures weight classes
equally, while micro-averaged measures combine their decisions. For a
multi-label task, threshold-dependent and ranking metrics can move in
different directions. This thesis fixes its principal measures before
comparison and preserves per-class outcomes, rather than changing the main
criterion after observing favourable results.

Qualitative evidence is useful when its meaning is precise. A partition
shows membership, a graph plot shows selected connectivity, and a weight
map shows a particular aggregation mechanism. None automatically measures
semantic region quality. In particular, similarity-based feature edges
can connect predictive contexts without representing named relationships
between objects. Evaluating graph reasoning therefore benefits from showing
both corrected examples and newly introduced errors, together with full-set
classification results. Selected successful images alone cannot establish
a general explanation of model behaviour.

Efficiency likewise belongs to the research question. Pooling reduces the
number of nodes in a graph but does not eliminate backbone computation or
the cost of constructing the graph. Parameter count, forward latency,
throughput, memory and total training duration describe different aspects
of that cost. Comparisons should state whether inputs are already on the
device and whether loading, preprocessing and transfer are included. An
added branch can have modest memory overhead yet affect latency through
its sequence of operations; this must be measured under shared conditions.

\section{Synthesis and research position}

The literature supports three distinctions that guide the experimental
questions in Chapter~\ref{chap:introduction}. First, a strong contextual
backbone and an explicit regional graph are different forms of organisation.
Their combination is plausible, but its incremental benefit must be
established through a common-backbone comparison. Second, learnable pooling
weights, learnable membership and learned graph attention are separate
sources of flexibility. Increasing one does not establish that the resulting
regions better represent land-cover entities or improve scene recognition.
Third, aggregation and feedback have clear precedents, so the present work
is positioned around a defined regional construction and its evaluation.

GAViT accordingly combines a pretrained Swin-T hierarchy with anchored region
aggregation, a sparse spatial-and-feature graph, graph attention and feedback
to the original tokens. Its fixed and local-prototype variants share the
main interfaces while changing the membership mechanism. Chapter~\ref{chap:method}
specifies the implementation; Chapter~\ref{chap:setup} separates the data
and model-selection protocols needed to assess it.

The intended evidence includes classification performance across tasks and
training repetitions, component comparisons, computational cost and
diagnostic visualisations. These support a bounded account of when the
added structure helps, when it harms, and which explanations remain
unresolved. This positioning allows positive and negative observations to
inform the same research questions without presupposing that a more
elaborate architecture must outperform its shared backbone.
